\chapter{Orientasi: Platform, API, dan Capstone}

\section*{Tujuan Pembelajaran}
\begin{itemize}
  \item Menjelaskan perbedaan bisnis \textit{pipeline} dan bisnis platform dengan contoh dari sektor yang dikenal.
  \item Menjelaskan cara kerja API (\textit{Application Programming Interface}) dengan bahasa sehari-hari, termasuk \textit{request}, \textit{response}, dan data JSON (\textit{JavaScript Object Notation}).
  \item Mengidentifikasi pertukaran data utama antara organisasi capstone dan pihak luarnya.
\end{itemize}

\section{Pendahuluan}

% TODO
% Cakupan: organisasi bertukar data dengan banyak pihak luar lewat surel, berkas, portal, dan
% sambungan sistem yang berbeda-beda; akibatnya antrean, kerja ulang, dan biaya integrasi naik.

\section{Peta Mata Kuliah dan Capstone Bersama}

% TODO
% Cakupan: alur 14 pertemuan dalam empat Bagian; tonggak C1 sampai C4 dan luaran mata kuliah ini
% di tiap tonggak; pembagian materi dengan IT40513, IT40713, dan IT40813; cara belajar mandiri
% dengan modul dan slide, serta fungsi pertemuan Sabtu untuk membahas progres.
% Kasus Medikara: posisi Medikara sebagai kasus cadangan bagi tim tanpa akses organisasi.

\section{Bisnis Pipeline dan Bisnis Platform}

% TODO
% Cakupan: rantai nilai linier yang membuat lalu menjual; platform yang mempertemukan pihak dan
% memfasilitasi interaksi; contoh lokal (transportasi daring, marketplace, dompet digital,
% telekonsultasi); tidak semua organisasi perlu menjadi platform.
% Kasus Medikara: rumah sakit sebagai pipeline dan Medikara+ sebagai calon platform.

\section{API dalam Bahasa Sehari-hari}

% TODO
% Cakupan: analogi loket layanan antarsistem; request, response, dan endpoint; JSON sebagai format
% data yang dapat dibaca manusia; kode status HTTP (HyperText Transfer Protocol) secara singkat.
% Kasus Medikara: aplikasi Medikara+ meminta jadwal dokter ke SIMRS.
% Pendalaman teknis: memanggil satu API publik tanpa akun dengan Postman atau curl, lalu membaca
% kode status dan isi JSON.

\section{Cara Organisasi Bertukar Data}

% TODO
% Cakupan: pertukaran manual lewat surel dan telepon; unggah berkas CSV atau spreadsheet; portal
% web; API; kelebihan dan risiko tiap cara; gejala awal integrasi point-to-point.
% Kasus Medikara: laporan manajemen berbasis spreadsheet, integrasi point-to-point ke BPJS
% Kesehatan dan SATUSEHAT.

\section{Kosakata Bersama untuk Tim Campuran}

% TODO
% Cakupan: istilah inti platform (sisi pasar, efek jaringan, ekosistem, mitra); istilah inti API
% (kontrak, gateway, token, versi); cara anggota teknis menjelaskan istilah ke anggota non-teknis
% dan sebaliknya; rujukan ke Daftar Singkatan.

\section{Kasus Berjalan Medikara Health}

% TODO
% Cakupan: profil organisasi fiktif PT Medikara Nusa Sehat; kondisi saat ini; strategi 2026 sampai
% 2030 "Satu Pasien, Satu Rekam Medis"; pemicu perubahan; etika memakai data organisasi nyata
% (izin dan anonimisasi).
% Kasus Medikara: seluruh ringkasan kasus, ditegaskan sebagai organisasi fiktif.

\section{Aktivitas: Membentuk Tim dan Memetakan Pertukaran Data}

% TODO
% Tujuan aktivitas: tim terbentuk, organisasi capstone terpilih, dan pertukaran data dengan pihak
% luar terdaftar sebagai bahan awal C1.
% Peran dalam tim: anggota berlatar bisnis, hukum, atau kesehatan memimpin wawancara singkat dan
% penelusuran dokumen tentang pihak luar; anggota berlatar teknis mencatat cara pertukaran data
% (surel, berkas, portal, API).
% Langkah: bentuk tim campuran; pilih organisasi atau Medikara; isi profil organisasi satu halaman;
% daftar pihak luar; catat data yang dipertukarkan, cara, frekuensi, dan masalah yang dirasakan.
% Luaran dan tonggak capstone: profil organisasi dan daftar pertukaran data, masuk ke C1.
% Berkas sources yang direncanakan: sources/chapter01/profil-organisasi.md,
% sources/chapter01/daftar-pertukaran-data.md, sources/chapter01/panduan-wawancara-singkat.md

\section{Ringkasan}

% TODO

\section{Latihan}

\begin{enumerate}
  \item Pilih satu layanan digital yang sering dipakai. Tentukan apakah layanan itu berpola \textit{pipeline} atau platform, lalu beri dua alasan.
  \item Jelaskan kepada rekan non-teknis apa yang terjadi ketika aplikasi pasien menampilkan jadwal dokter, dengan memakai istilah \textit{request}, \textit{response}, dan API.
  \item Bandingkan pengiriman data klaim asuransi lewat unggah berkas dan lewat API. Sebutkan satu kelebihan dan satu risiko masing-masing.
\end{enumerate}
